\documentclass[11pt]{article}
\usepackage[margin=1in]{geometry}
\usepackage{amsmath,amssymb,amsthm,booktabs,array,enumitem,hyperref}
\newtheorem{theorem}{Theorem}
\newtheorem{proposition}[theorem]{Proposition}
\newtheorem{lemma}[theorem]{Lemma}
\theoremstyle{definition}
\newtheorem{definition}[theorem]{Definition}
\newcommand{\CVaR}{\operatorname{CVaR}}
\newcommand{\argmax}{\operatorname*{arg\,max}}
\title{Canonical Repaired Theorem Set\\Multi-Crop Ranking Reversal}
\author{Issue \#2 theory audit deliverable}
\date{19 July 2026}

\begin{document}
\maketitle

\section{Canonical model}

On $(\Omega,\mathcal F,\mathbb P)$ let $\mathcal I=\{1,\ldots,n\}$ index crops.
The ex-ante acreage decision is $x\in\mathbb R_+^n$.  Per-hectare crop
profit is $\pi_i=P_iY_i-c_i$, portfolio profit is
$\Pi(x)=\pi^\top x$, and loss is
\[
L(x)=-\Pi(x),\qquad \mu=\mathbb E[\pi].
\]
The retained operational set is
\[
X=\{x:\mathbf1^\top x\le A,\ a^\top x\le B,\ \ell\le x\le u,\
Gx\le g\}.
\]
Assume throughout that $X$ is nonempty and compact and every $\pi_i$ is
integrable.  For $\alpha\in(0,1)$ define loss CVaR by
\[
r_\alpha(x)=\inf_{v\in\mathbb R}
\left\{v+\frac{1}{1-\alpha}\mathbb E[(L(x)-v)_+]\right\}.
\]
The canonical program and solution correspondence are
\[
(P_\kappa)\quad
\max_{x\in X}\{\mu^\top x:r_\alpha(x)\le\kappa\},\qquad
S(\kappa)=\argmax(P_\kappa).
\]
Marginal laws may be joined by a valid $n$-copula $C_\theta$.
The bivariate lower-tail coefficient, when its limit exists, is
$\lambda_L=\lim_{q\downarrow0}C(q,q)/q$; it is a diagnostic, not an
ordering of general joint laws.

\section{Reversal definitions}

Let $s_i>s_j$ and let $\tau_x>0$ be a declared numerical tolerance.
A pair $(i,j)$ is a possible reversal if some $x\in S(\kappa)$ has
$x_j\ge x_i+\tau_x$, a universal reversal if every such $x$ does, and
a selected reversal if a declared selection $x^\sigma\in S(\kappa)$
does.  A strong reversal has $x_i=0<x_j$.  A top-rank reversal compares
a highest-score crop with a lower-score crop.  These labels are not
interchangeable.

\section{Results and proofs}

\begin{theorem}[CT1: existence, convexity, and finite LP]\label{ct1}
$r_\alpha$ is finite, continuous, and convex on $X$.  If
$(P_\kappa)$ is feasible, $S(\kappa)$ is nonempty and compact.  Under
finite scenarios $\pi_s$ with weights $w_s>0$, $\sum_s w_s=1$, its risk
constraint is represented exactly by
\[
v+\frac{1}{1-\alpha}\sum_s w_sq_s\le\kappa,\quad
q_s\ge-\pi_s^\top x-v,\quad q_s\ge0,\quad v\in\mathbb R.
\]
\end{theorem}
\begin{proof}
Each $x\mapsto-\pi(\omega)^\top x$ is affine.  The positive part and
expectation preserve convexity, and minimization in the Rockafellar--Uryasev
representation gives the convex CVaR functional.  Integrability and compact
$X$ give finiteness and continuity.  Its sublevel set intersected with $X$
is compact; a continuous linear objective attains its maximum there.
For finite support, replacing every positive part by its epigraph variable
$q_s$ gives the displayed inequalities, and minimizing over $q_s$ and free
$v$ recovers the definition exactly.
\end{proof}

\begin{theorem}[CT2: ordinal ranking insufficiency]\label{ct2}
A strict weak order induced by suitability scores does not identify a
cardinal acreage vector or an optimal allocation without an additional
mapping from scores to payoffs and constraints.
\end{theorem}
\begin{proof}
Every strictly increasing transformation of the scores preserves their
order.  Infinitely many cardinal vectors and feasible sets are compatible
with that same order, and any selected feasible vector can be made uniquely
optimal by a suitable linear objective over a suitable compact polytope.
Thus the ordinal information alone cannot identify acreage.
\end{proof}

\begin{proposition}[CT3: feasibility-forced universal reversal]\label{ct3}
If $s_i>s_j$ and
\[
\sup_{x\in X}x_i+\tau_x\le\inf_{x\in X}x_j,
\]
then every feasible point, hence every optimizer, reverses $(i,j)$.
The readily checked condition $u_i+\tau_x\le\ell_j$ is sufficient.
\end{proposition}
\begin{proof}
For every $x\in X$, $x_i+\tau_x\le\sup_Xx_i+\tau_x
\le\inf_Xx_j\le x_j$.  Restriction to $S(\kappa)\subseteq X$ preserves
the inequality.
\end{proof}

\begin{theorem}[CT4: exact risk-slack identity]\label{ct4}
Let $S_0=\argmax_{x\in X}\mu^\top x$.  Then
\[
S(\kappa)=S_0\cap\{x:r_\alpha(x)\le\kappa\}
\]
whenever the right side is nonempty.  Consequently all unconstrained
optima remain optimal if all are risk-feasible; a single slack optimum
does not imply equality of optimal sets under multiplicity.
\end{theorem}
\begin{proof}
No risk-feasible point can exceed the unrestricted optimum value.
Every risk-feasible member of $S_0$ attains that value, while any
risk-constrained optimizer attaining a smaller value would be dominated by
such a member.  Hence the two sets coincide.
\end{proof}

\begin{theorem}[CT5: complete KKT and finite tail weights]\label{ct5}
Assume a convex constraint qualification.  A feasible $x^*$ is optimal
if and only if there exist nonnegative multipliers
$\eta,\beta,\rho,\gamma,\underline\nu,\overline\nu$ and
$d\in\partial r_\alpha(x^*)$ satisfying
\[
0=-\mu+\eta\mathbf1+\beta a+G^\top\gamma
-\underline\nu+\overline\nu+\rho d
\]
together with primal feasibility and complementary slackness for land,
budget, every row of $G$, both bound systems, and risk.  Under finite
scenarios one may take
\[
d=-\sum_s\xi_s\pi_s,\quad
0\le\xi_s\le\frac{w_s}{1-\alpha},\quad \sum_s\xi_s=1,
\]
with upper-tail saturation and VaR-atom splitting dictated by LP
complementarity.
\end{theorem}
\begin{proof}
Apply convex KKT to minimization of $-\mu^\top x$ with all inequalities
written in nonpositive form.  The displayed stationarity is the
subdifferential inclusion after collecting their normals; KKT necessity
follows from the constraint qualification and sufficiency from convexity.
For finite support, dualizing the excess-loss epigraph bounds yields the
box constraints on $\xi_s$, while stationarity in the free variable $v$
gives $\sum_s\xi_s=1$.  Stationarity in $x$ yields the displayed
subgradient, and complementarity assigns full weight above VaR, zero below,
and admissible fractional weight at atoms.
\end{proof}

\begin{proposition}[CT6: risk-adjusted feasible displacement]\label{ct6}
Let $x$ be feasible and let $h\ne0$ satisfy $x+th\in X$ for all small
$t\ge0$.  Suppose $\mu^\top h>0$ and either (i) $r_\alpha(x)<\kappa$,
(ii) $r_\alpha(x)=\kappa$ and $r_\alpha'(x;h)<0$, or (iii) direct
evaluation establishes $r_\alpha(x+th)\le\kappa$ for some small $t>0$.
Then $x$ is not optimal.
\end{proposition}
\begin{proof}
In case (i), continuity makes every sufficiently small step risk-feasible.
In case (ii), directional differentiability gives
$r_\alpha(x+th)=r_\alpha(x)+t r_\alpha'(x;h)+o(t)<\kappa$ for small
positive $t$.  Case (iii) supplies a feasible step directly.  In all cases
the objective rises by $t\mu^\top h>0$.  Strict negativity in (ii) cannot
be weakened to zero: an active boundary locally shaped as $t^2\le0$ has
zero derivative but permits no positive step.
\end{proof}

\begin{theorem}[CT7: restricted dependence ordering]\label{ct7}
Fix the marginals and a named copula family.  If for
$\theta_1\preceq\theta_2$ and every $x$ in a declared domain $D\subseteq X$
\[
L_{\theta_1}(x)\preceq_{\rm cx}L_{\theta_2}(x),
\]
then $r_{\alpha,\theta_1}(x)\le r_{\alpha,\theta_2}(x)$ on $D$.
Consequently the higher-dependence risk-feasible set within $D$ is a
subset of the lower-dependence one.
\end{theorem}
\begin{proof}
For each $v$, $z\mapsto(z-v)_+$ is convex.  Convex order therefore orders
its expectations.  Add $v$, scale by $(1-\alpha)^{-1}$, and take the
infimum over $v$ to obtain the CVaR inequality.  Sublevel-set inclusion is
immediate.  No claim follows from $\lambda_L$ alone.
\end{proof}

\begin{theorem}[CT8: reversal regions and crossing sets]\label{ct8}
For a declared selection rule $\sigma$, define
$\Delta^\sigma_{ij}(\theta)=x^\sigma_j(\theta)-x^\sigma_i(\theta)$ and
\[
\mathcal R^\sigma_{ij}=\{\theta:\Delta^\sigma_{ij}(\theta)\ge\tau_x\},
\qquad
\mathcal C^\sigma_{ij}=\{\theta:\Delta^\sigma_{ij}(\theta)=0\}.
\]
Without continuity, uniqueness, strict monotonicity, endpoint crossing,
and active-set regularity, the crossing set may be empty, a singleton,
multiple points, or an interval.
\end{theorem}
\begin{proof}
The set definitions are exact.  Constant positive and constant negative
gaps give empty crossing sets; an affine sign-changing gap gives one point;
a polynomial with several roots gives several; and a function vanishing on
an interval gives an interval.  Optimization correspondences can realize
discontinuity and multiplicity, so none of the omitted properties follows
from the canonical program.
\end{proof}

\begin{proposition}[CT9: pseudo-diversification is diagnostic]\label{ct9}
For predeclared thresholds $\rho_0,\lambda_0$, the label
$D_{\rho,\lambda}=\{(i,j):|\rho_{ij}|\le\rho_0,\
\lambda_{L,ij}\ge\lambda_0\}$ is descriptive only.  Membership implies
neither optimal inclusion or exclusion nor an ordering of portfolio CVaR
or welfare.
\end{proposition}
\begin{proof}
Optimal inclusion also depends on means, all constraints, the full joint
law, and other crops.  Raising a flagged crop's mean sufficiently can force
inclusion; an upper bound or low mean can force exclusion independently of
the flag.  Equal scalar diagnostics can coexist with distinct joint laws
and CVaR.  Therefore no asserted optimization conclusion is logically
entailed.
\end{proof}

\begin{theorem}[CT10: information actionability and weak flexibility]\label{ct10}
Let a signal $Z$ arrive before a recourse action is selected.  If ignoring
$Z$ is admissible, posterior optimization has value no lower than prior
optimization.  Its value of information is zero when one common policy is
optimal for almost every posterior problem.  If
$\mathcal A_{f_1}(z)\subseteq\mathcal A_{f_2}(z)$ almost surely, optimized
value is weakly increasing in flexibility.  Strict increase and
supermodular information--flexibility complementarity do not follow
without additional structure.
\end{theorem}
\begin{proof}
The posterior decision maker can copy the prior policy, so optimization
over signal-contingent policies weakly improves value.  A common posterior
optimizer attains both values, proving zero value.  Maximization over a
superset cannot lower value, proving weak flexibility monotonicity.
Equality occurs whenever the added actions are never optimal, refuting
strictness; substitutes between information and recourse refute general
supermodularity.
\end{proof}

\section{Transition and evidence boundary}

The Draft's multi-crop architecture is retained.  The false general
ranking-reversal equivalence is replaced by CT3, CT5, and CT6; global
tail-dependence monotonicity by conditional CT7; a unique threshold by
CT8; the pseudo-diversification proposition by CT9; and strict
information--flexibility complementarity by CT10.  Every Draft result
R01--R31 is dispositioned in the transition registry.  All teacher-Draft
numbers remain illustrative and are inadmissible as paper evidence until
their separate simulation or empirical contracts are executed.

\section*{Method sources}

The loss-CVaR representation and finite-scenario LP follow the full-text
method foundations of Rockafellar and Uryasev (2000, 2002).  Dependence
language is restricted to named, verified families, consistent with the
family-specific analysis of Ansari and Rockel (2024).  Exact source URLs,
DOIs, verification dates, roles, and limitations are recorded in the
literature registry and method-source verification note.

\end{document}
